\paragraph{El producto cruz componente por componente.} Los lados $\vect{e}_1$ y $\vect{e}_2$ del triángulo anterior son vectores, y su producto cruz $\vect{n}=\vect{e}_1\times\vect{e}_2$ es perpendicular al triángulo. Sus componentes (figura~\ref{fig:vec-triangulo}) salen así:
\begin{align*}
n_x&=e_{1y}e_{2z}-e_{1z}e_{2y}=\left(-1{,}1582\right)\cdot0{,}1548-\left(-0{,}2729\right)\cdot\left(-0{,}0189\right)=-0{,}18445,\\
n_y&=e_{1z}e_{2x}-e_{1x}e_{2z}=\left(-0{,}2729\right)\cdot1{,}1917-\left(-0{,}0033\right)\cdot0{,}1548=-0{,}32470,\\
n_z&=e_{1x}e_{2y}-e_{1y}e_{2x}=\left(-0{,}0033\right)\cdot\left(-0{,}0189\right)-\left(-1{,}1582\right)\cdot1{,}1917=1{,}38029,
\end{align*}
es decir $\vect{n}=\begin{pmatrix} -0{,}1844 \\ -0{,}3247 \\ 1{,}3803 \end{pmatrix}$. Su norma es $\lVert\vect{n}\rVert=1{,}4299$, que es el doble del área del triángulo (área $=0{,}7150$ mm$^{2}$), y
el vector unitario es $\hat{\vect{n}}=\vect{n}/\lVert\vect{n}\rVert=\begin{pmatrix} -0{,}1290 \\ -0{,}2271 \\ 0{,}9653 \end{pmatrix}$. Como $\vect{n}$ es perpendicular al triángulo, su producto punto
con los dos lados debe ser cero: $\hat{\vect{n}}\cdot\vect{e}_1=0{,}0\times10^{0}$ y $\hat{\vect{n}}\cdot\vect{e}_2=-2{,}8\times10^{-17}$ (los ceros de la
máquina vienen de redondear a cuatro decimales). Y el producto punto de la normal con el rayo hacia el centro del triángulo,
$\vect{c}_t=\vect{p}_0+\vect{p}_1+\vect{p}_2=\begin{pmatrix} 0{,}536 \\ 0{,}044 \\ 241{,}535 \end{pmatrix}$, dice si la cámara lo ve de frente o de canto:
$\lvert\hat{\vect{n}}\cdot\hat{\vect{c}}_t\rvert=0{,}9650$ (un valor cercano a $1$ es de frente y cercano a $0$, de canto).

\begin{figure}[htbp]
\centering
\includegraphics[width=0.97\linewidth]{vec_triangulo.jpg}
\caption{Triángulo de la malla del casco: los lados $\vect{e}_1$ y $\vect{e}_2$, la normal unitaria $\hat{\vect{n}}$ del producto cruz y
las operaciones con vectores (norma, producto punto con los lados y área).}
\label{fig:vec-triangulo}
\end{figure}
